\documentclass{beamer}

\usetheme{metropolis}
\usepackage[french]{babel}
\usepackage[utf8]{inputenc}
\usepackage[T1]{fontenc}
\usepackage{graphicx}
\usepackage{tikz}
\newcommand{\roundimage}[2][2cm]{%
  \begin{tikzpicture}
    \clip[rounded corners=1cm] (0,0) rectangle (#1,#1);
    \node[anchor=south west, inner sep=0] at (0,0) {\includegraphics[width=#1]{#2}};
  \end{tikzpicture}%
}
\usepackage{pifont} % À mettre dans le préambule

% Icônes utiles
\newcommand{\question}{\ding{72}}   % ❓
\newcommand{\blocked}{\ding{55}}    % 🚫
\newcommand{\lock}{\ding{115}}      % 🔐
\newcommand{\target}{\ding{93}}     % 🎯


\setbeamertemplate{title page}{
    \vbox{}
    \vfill
    \begin{center}
        {\small % ⬅️ Réduction ici
        
        % Logos en haut
        \begin{minipage}{0.45\textwidth}
            \flushleft
            \includegraphics[height=1.2cm]{logo}
        \end{minipage}%

        \vspace{0.5cm}

        {\usebeamerfont{title}\Huge \inserttitle\par}
        \vspace{0.3cm}
        {\usebeamerfont{subtitle}\Large \insertsubtitle\par}
        \vspace{1cm}
        {\usebeamerfont{institute}\normalsize \insertinstitute\par \usebeamerfont{date}\insertdate\par}
        
        } % fin du \small
    \end{center}
    \vfill
}

\setbeamertemplate{section in toc}[sections numbered]


% Informations
\title[Soutenance de Projet]{Soutenance de Projet}
\subtitle{SAE-821 Développement d’un
moteur de recommandations
respectueux de la vie privée}
\author{
    \textbf{Enzo Berge, Lois Hermelin}\\
    \textbf{Florian Audouard, Mattéo Pegliasco}
}
\institute{Université de Toulon, La Garde}
\date{2 Juin 2025}

\begin{document}

% Page de titre
\begin{frame}
    \titlepage
\end{frame}

% Sommaire
\begin{frame}
    \frametitle{Sommaire}
    \tableofcontents
\end{frame}

% Introduction
\section{Contexte}
\begin{frame}
    \frametitle{Contexte}
    \begin{itemize}
        \item Les moteurs de recommandation sont omniprésents : plateformes de streaming, e-commerce, réseaux sociaux.
        \item Leur objectif est de proposer du contenu pertinent basé sur les préférences, l’historique ou les comportements similaires.
        \item Pour cela, ils exploitent d’importantes quantités de données personnelles.
        \item Cette collecte massive soulève des enjeux cruciaux en matière de confidentialité.
        \item Dans ce projet, nous cherchons à concevoir un moteur de recommandations simple mais performant, respectueux de la vie privée.
    \end{itemize}
\end{frame}


% Problématique
\section{Problématique}
\begin{frame}
    \frametitle{Problématique}

    \small % Réduit la taille du texte

    \textbf{\question~Problème central :} 
    \textit{Comment proposer des recommandations pertinentes sans compromettre la vie privée des utilisateurs ?}

    \vspace{0.2cm}
    \textbf{\blocked~Limites actuelles :}
    \begin{itemize}
        \item Solutions \textbf{centralisées} : profilage, tracking, fuites de données.
        \item Collecte massive de données personnelles.
    \end{itemize}

    \textbf{\lock~Deux approches de protection :}
    \begin{itemize}
        \item \textbf{Differential Privacy} : bruit pour anonymiser.
        \item \textbf{Cryptographie homomorphe} : calcul sur données chiffrées.
    \end{itemize}

    \textbf{\target~Objectif :} 
    Développer un moteur de recommandation basé sur \texttt{MovieLens}, tout en garantissant la confidentialité.
\end{frame}

% User Story
\section{User Story}
\begin{frame}
    \frametitle{User Story}
    \begin{block}{Exemple}
        \textit{“En tant qu’utilisateur, je veux recevoir des recommandations personnalisées sans que mes données personnelles soient stockées en ligne.”}
    \end{block}
    \begin{itemize}
        \item Vision utilisateur claire
        \item Respect des principes de confidentialité dès la conception
    \end{itemize}
\end{frame}

% Fonctionnalités - Vue générale
\section{Fonctionnalités}
\begin{frame}
    \frametitle{Fonctionnalités - Vue générale}
    \begin{itemize}
        \item Recommandations pertinentes
        \item Pas de collecte centralisée
        \item Interface intuitive
        \item Système modulaire
    \end{itemize}
\end{frame}

% Fonctionnalités - Techniques
\begin{frame}
    \frametitle{Fonctionnalités - Techniques}
    \begin{itemize}
        \item Algorithmes de filtrage collaboratif
        \item Traitement local des données utilisateur
        \item Anonymisation des échanges entre profils
        \item Visualisation des recommandations avec Streamlit
    \end{itemize}
\end{frame}

% Architecture - Diagramme UML
\section{Architecture}
\begin{frame}
    \frametitle{Architecture du système}
    \begin{center}
        \includegraphics[width=0.8\linewidth]{logo}
        % Remplacez "uml_diagram" par le nom réel de votre image UML
    \end{center}
\end{frame}

% Architecture - Explication
\begin{frame}
    \frametitle{Modules principaux}
    \begin{itemize}
        \item \textbf{Interface utilisateur} : Streamlit pour l’interaction
        \item \textbf{Moteur de recommandation} : filtrage collaboratif local
        \item \textbf{Base de données locale} : stockage des interactions
        \item \textbf{Système d’anonymisation} : confidentialité des échanges
    \end{itemize}
\end{frame}

% Déroulement
\section{Déroulement du projet}
\begin{frame}
    \frametitle{Déroulement du projet}
    \begin{itemize}
        \item Méthodologie agile en sprint de 2 semaines
        \item Outils utilisés :
        \begin{itemize}
            \item Trello : gestion des tâches
            \item GitHub : collaboration, versioning
            \item Streamlit : front-end
        \end{itemize}
        \item Répartition claire des rôles
    \end{itemize}
\end{frame}

% Résultats
\section{Résultats}
\begin{frame}
    \frametitle{Résultats obtenus}
    \begin{itemize}
        \item Prototype fonctionnel
        \item Tests de recommandation concluants
        \item Interface ergonomique
        \item Données utilisateurs protégées
    \end{itemize}
\end{frame}

% Conclusion
\section{Conclusion}
\begin{frame}
    \frametitle{Conclusion}
    \begin{itemize}
        \item Projet mené avec rigueur et efficacité
        \item Objectifs respectés
        \item Confidentialité au cœur du développement
        \item Enrichissement personnel et technique
    \end{itemize}
\end{frame}

% Perspectives
\section{Perspectives}
\begin{frame}
    \frametitle{Perspectives}
    \begin{itemize}
        \item Intégration d’autres types de contenus : vidéos, podcasts
        \item Optimisation de l’algorithme (auto-apprentissage)
        \item Tests utilisateurs pour améliorer l’expérience
        \item Déploiement à plus grande échelle
    \end{itemize}
\end{frame}

\begin{frame}[c]
    \centering
    \includegraphics[width=0.4\textwidth]{logo.png}

    \vspace{1cm}
    {\Huge \textbf{Merci pour votre attention !}}

    \vspace{1.5cm}
    {\Large N'hésitez pas à poser vos questions !}
\end{frame}

\end{document}